\documentclass[11pt]{article}
\usepackage[T1]{fontenc}
\usepackage[utf8]{inputenc}
\usepackage{lmodern}
\usepackage[margin=1in]{geometry}
\usepackage{amsmath,amssymb,amsthm,mathtools}
\usepackage{microtype,booktabs,array}
\usepackage[hidelinks]{hyperref}
\usepackage{fancyvrb}
\hypersetup{pdftitle={A rational completely positive matrix with no rational nonnegative factor},pdfsubject={PF-03: exact counterexample and unrestricted-width obstruction},pdfauthor={OpenAI}}
\setlength{\parskip}{0.35em}
\setlength{\parindent}{1.1em}
\setlength{\emergencystretch}{2em}
\newtheorem{theorem}{Theorem}[section]
\newtheorem{lemma}[theorem]{Lemma}
\newtheorem{proposition}[theorem]{Proposition}
\newtheorem{corollary}[theorem]{Corollary}
\theoremstyle{definition}
\newtheorem{definition}[theorem]{Definition}
\newtheorem{remark}[theorem]{Remark}
\newcommand{\R}{\mathbb R}
\newcommand{\Q}{\mathbb Q}
\newcommand{\Z}{\mathbb Z}
\newcommand{\F}{\mathbb F}
\newcommand{\CP}{\mathcal{CP}}
\newcommand{\cone}{\operatorname{cone}}
\newcommand{\rank}{\operatorname{rank}}
\newcommand{\tr}{\operatorname{tr}}
\newcommand{\range}{\operatorname{range}}
\newcommand{\cpr}{\operatorname{cpr}}
\newcommand{\spanof}{\operatorname{span}}
\newcommand{\ri}{\operatorname{ri}}
\newcommand{\trans}{^{\mathsf T}}
\numberwithin{equation}{section}

\begin{document}
\begin{center}
{\LARGE\bfseries A rational completely positive matrix\\[0.2em]
with no rational nonnegative factor}
\end{center}

\begin{abstract}
An explicit integer matrix of order $444$ and rank seven is constructed that
is completely positive over the reals but has no nonnegative rational Gram
factor of any finite width. All its entries are strictly positive, and its real
cp-rank is seven. The construction gives a negative answer to the universal
boundary-factorization question PF-03. The obstruction is a rational
polyhedral cone in $\R^7$ on which a trace-zero quadratic form is nonnegative
and has exactly seven irrational zero rays forming an orthonormal basis.
The arithmetic seed lies in $\Q(\sqrt[3]{2})$. Rational interval calculations
certify its signs, and exhaustive integer computation certifies all $444$
facets of the cone. The complete matrix, a compact integer representation,
and a standard-library Python verifier accompany the proof. This is a
computer-assisted proof presented for independent review; it has not been
formalized in a proof assistant or independently reviewed.
\end{abstract}

\section{The statement being refuted}
A real symmetric matrix $A$ is completely positive if $A=BB\trans$ for some
entrywise nonnegative real matrix $B$ with finitely many columns. Write
$\CP_n$ for the cone of such matrices of order $n$. The question in PF-03 is
whether every
\begin{equation}\label{eq:question}
 A\in\Q^{n\times n}\cap\partial\CP_n,\qquad n\geq5,
\end{equation}
admits $A=CC\trans$ with $C\in\Q_{\geq0}^{n\times m}$ for some finite $m$.
The width is explicitly unrestricted. Boundaries are in the full space of
real symmetric matrices.\footnote{See the exact PF-03 statement in
\cite{repo}, especially its unrestricted-width qualification.}

The distinction between a rational Gram factor and a \emph{nonnegative}
rational Gram factor is essential. Likewise, failure of a square rational
factor does not suffice: a wider rational factor could still exist. The
argument below excludes every finite width. The interior theorem and the
boundary question are stated separately in the 2025 survey of Berman and
Shaked-Monderer.\footnote{\cite{bsm}, Section 2, Question 2.1, Theorem 2.2,
and the boundary Open Problem on page 2. The interior theorem is also
\cite{dsv}, Theorem 1.1.}

\begin{theorem}[Explicit counterexample]\label{thm:main}
Let $R\in\Z^{444\times7}$ be the primitive integer facet matrix defined in
Sections~\ref{sec:cone}--\ref{sec:facets}, in the row order specified there,
and set
\begin{equation}\label{eq:Adefinition}
 A=RR\trans.
\end{equation}
Then all entries of $A$ are strictly positive integers,
\[
 \rank(A)=\cpr(A)=7,\qquad A\in\partial\CP_{444},
\]
and there is no $C\in\Q_{\geq0}^{444\times m}$ satisfying $CC\trans=A$,
for any finite integer $m\geq1$.
\end{theorem}

The exact $444\times7$ array is supplied as \texttt{data/R\_integer.csv}.
The expanded lower triangle of $A$ is supplied as
\texttt{data/A\_integer.mtx.gz}. Thus the matrix is available both explicitly
and in a compact exact representation. Equation~\eqref{eq:Adefinition} is
\emph{not} its nonnegative factorization: $R$ is a signed integer matrix.
Its nonnegative real factor is $RO$, where $O$ is constructed below.

\section{A general unrestricted-width obstruction}\label{sec:obstruction}
For a cone $K\subset\R^r$, write
\[
 \CP(K)=\cone\{xx\trans:x\in K\}.
\]
All conic combinations here are finite. The following lemma isolates the
part of the argument that excludes arbitrary factor widths.

\begin{lemma}[Cone obstruction]\label{lem:obstruction}
Suppose that $K\subset\R^r$ is a full-dimensional pointed rational
polyhedral cone and
\[
 K=\{x\in\R^r:Rx\geq0\},\qquad R\in\Q^{N\times r},\quad \rank R=r.
\]
Suppose $O=[u_1\ \cdots\ u_r]\in\R^{r\times r}$ is orthogonal and every
$u_i\in K$. Suppose also that a real symmetric matrix $Q$ satisfies
\begin{equation}\label{eq:abstractconditions}
 \tr Q=0,\qquad x\trans Qx\geq0\ (x\in K),\qquad
 \{x\in K\cap\Q^r:x\trans Qx=0\}=\{0\}.
\end{equation}
Then $A=RR\trans$ is completely positive and admits no nonnegative rational
Gram factor of any finite width. If $N>r$, it is a boundary matrix of $\CP_N$.
\end{lemma}

\begin{proof}
Since $RO\geq0$ and $OO\trans=I_r$,
\begin{equation}\label{eq:realCP}
 A=(RO)(RO)\trans\in\CP_N.
\end{equation}
Assume, for a contradiction, that $A=CC\trans$ with a finite nonnegative
rational matrix $C=[c_1\ \cdots\ c_m]$.

Every column of $C$ lies in $\range R$. Indeed, if $y\in\ker R\trans$, then
\[
 0=y\trans Ay=\sum_{j=1}^m(y\trans c_j)^2.
\]
All summands are nonnegative real numbers, so $y\trans c_j=0$ for every $j$.
Taking the orthogonal complement of $\ker R\trans$ gives the assertion.

The matrix
\[
 L=(R\trans R)^{-1}R\trans
\]
is rational and satisfies $LR=I_r$. Put $X=LC=[x_1\ \cdots\ x_m]$.
Then $X$ is rational, $RX=C\geq0$, and
\begin{equation}\label{eq:reducedidentity}
 XX\trans=LAL\trans=I_r.
\end{equation}
In particular every $x_j\in K$. Now
\begin{equation}\label{eq:tracecontradiction}
 0=\tr Q=\tr(QXX\trans)=\sum_{j=1}^m x_j\trans Qx_j.
\end{equation}
Every summand is nonnegative. Each must vanish, and
\eqref{eq:abstractconditions} forces every $x_j=0$, contradicting
\eqref{eq:reducedidentity}. No upper bound or other restriction on $m$ has
been used.

If $N>r$, then $A$ is singular. For a nonzero $z\in\ker A$, the matrix
$A-\varepsilon zz\trans$ fails to be positive semidefinite for every
$\varepsilon>0$. Since every completely positive matrix is positive
semidefinite, no neighborhood of $A$ is contained in $\CP_N$. Together
with $A\in\CP_N$, this places $A$ on its boundary.
\end{proof}

We will certify a stronger zero-set property than the one in
\eqref{eq:abstractconditions}: the nonzero zeros in $K$ are exactly the
positive rays through the orthonormal vectors $u_1,\ldots,u_7$, and none of
these rays contains a nonzero rational vector.

\begin{lemma}[Irrational ray test]\label{lem:irrational}
Let $\alpha=\sqrt[3]{2}$ and
\[
 u=a+\alpha b+\alpha^2c,\qquad a,b,c\in\Q^r.
\]
If $\rank[a\ b\ c]>1$, then $\R u$ contains no nonzero rational vector.
\end{lemma}
\begin{proof}
The polynomial $t^3-2$ is irreducible over $\Q$ by Eisenstein's criterion,
so $1,\alpha,\alpha^2$ are linearly independent over $\Q$. Suppose
$u=sw$ with $0\ne w\in\Q^r$. Choose a nonzero coordinate $w_i$.
Then $s=u_i/w_i\in\Q(\alpha)$, so write
$s=s_0+s_1\alpha+s_2\alpha^2$ with $s_k\in\Q$.
Comparing coefficients gives $a=s_0w$, $b=s_1w$, and $c=s_2w$.
The coefficient matrix has rank at most one, a contradiction.
\end{proof}

\section{An exact seven-dimensional algebraic seed}\label{sec:seed}
Set $\alpha=\sqrt[3]{2}>0$ and $\F=\Q(\alpha)$. The field is viewed in
its unique real embedding. Define the following skew-symmetric integer
matrices:
\begin{equation}\label{eq:S}
 S=\begin{pmatrix}
 0&-1&0&-6&-2&-6&-4\\
 1&0&3&4&-1&5&-1\\
 0&-3&0&-4&1&-2&-5\\
 6&-4&4&0&1&3&-3\\
 2&1&-1&-1&0&-1&-1\\
 6&-5&2&-3&1&0&-3\\
 4&1&5&3&1&3&0
 \end{pmatrix},
\end{equation}
\begin{equation}\label{eq:T}
 T=\begin{pmatrix}
 0&0&-3&0&-2&2&5\\
 0&0&-4&2&2&-2&2\\
 3&4&0&1&1&-3&0\\
 0&-2&-1&0&4&3&0\\
 2&-2&-1&-4&0&-2&-5\\
 -2&2&3&-3&2&0&4\\
 -5&-2&0&0&5&-4&0
 \end{pmatrix}.
\end{equation}
Put
\begin{equation}\label{eq:cayley}
 \Omega=\alpha S+\alpha^2T,\qquad
 O=(I_7+\Omega)^{-1}(I_7-\Omega).
\end{equation}
The real skew-symmetry of $\Omega$ makes $I_7\pm\Omega$ invertible.
The Cayley transform is orthogonal, and all entries of $O$ belong to $\F$.
These identities are also checked directly in exact field arithmetic.

Expand $O$ uniquely as
\begin{equation}\label{eq:coefficientmatrices}
 O=B_0+\alpha B_1+\alpha^2B_2,\qquad B_k\in\Q^{7\times7}.
\end{equation}
For $1\leq i\leq7$, define
\begin{equation}\label{eq:Cis}
 C_i=[(B_0)_{:i}\ (B_1)_{:i}\ (B_2)_{:i}],\quad
 U_i=\range C_i,\quad v=(1,\alpha,\alpha^2)\trans,
 \quad u_i=C_iv=O_{:i}.
\end{equation}
The exact matrices $C_i$ all have rank three. Lemma~\ref{lem:irrational}
therefore applies to each $u_i$.

\subsection{Definition of the exposing quadratic form}
Let $Q'=(q_{ij})\in\F^{7\times7}$ be symmetric with zero diagonal. Its
$21$ upper-triangular entries are ordered lexicographically by $(i,j)$.
Fix the last seven entries by
\begin{equation}\label{eq:freevalues}
 \begin{array}{c|rrrrrrr}
 (i,j)&(3,7)&(4,5)&(4,6)&(4,7)&(5,6)&(5,7)&(6,7)\\ \hline
 1000q_{ij}&63&1007&18&402&151&1522&1640.
 \end{array}
\end{equation}
Define the remaining $14$ entries by the $14$ equations
\begin{equation}\label{eq:tangentlinear}
 \bigl(Q'O\trans B_k\bigr)_{ii}=0,
 \qquad 1\leq i\leq7,\quad k=0,1.
\end{equation}
The coefficient matrix has exact rank $14$, with pivots in the first $14$
lexicographic upper-triangular variables. Thus these equations and
\eqref{eq:freevalues} uniquely specify $Q'$. Finally set
\begin{equation}\label{eq:Q}
 Q=OQ'O\trans,\qquad H_i=C_i\trans QC_i.
\end{equation}

This defines the entire seed using only two small integer matrices, seven
rational numbers, and arithmetic in a specified cubic field. The files
\texttt{build\_seed.py} and \texttt{exact\_algebraic\_certificate.json}
provide, respectively, a deterministic reconstruction and all expanded
field elements. Each field element is stored as its three rational
coefficients in the basis $(1,\alpha,\alpha^2)$.

\begin{lemma}[Certified seed properties]\label{lem:seed}
For the preceding exact seed,
\begin{equation}\label{eq:seedproperties}
 OO\trans=O\trans O=I_7,\qquad \tr Q=0,
 \qquad q_{ij}>\frac{17391}{10^6}>0\quad(i\ne j).
\end{equation}
Every $C_i$ has rank three, and every $H_i$ is positive semidefinite of
rank two with kernel $\R v$.
\end{lemma}

\begin{proof}
Orthogonality follows from the Cayley transform; cyclic invariance of trace
and the zero diagonal of $Q'$ give $\tr Q=0$.
The exact coefficient computations give $\rank C_i=3$ for all seven
indices. Equations~\eqref{eq:tangentlinear}, together with
\[
 O\trans B_0+\alpha O\trans B_1+\alpha^2O\trans B_2=I_7,
\]
also imply $\bigl(Q'O\trans B_2\bigr)_{ii}=0$. Consequently
$C_i\trans Q u_i=0$, or equivalently $H_iv=0$.

The remaining assertions are finite sign checks in $\F$. Table~\ref{tab:local}
records strict rational lower bounds for the two leading principal minors
needed for each $H_i$. The bound in \eqref{eq:seedproperties} is also obtained
by exact rational interval evaluation. The sign algorithm and isolating
interval are specified in Section~\ref{sec:verification}; they use no
floating-point tolerance.

For completeness, positivity of those two minors and $H_iv=0$ really do
prove the asserted semidefiniteness. Write
\[
 H_i=\begin{pmatrix}D&b\\b\trans&c\end{pmatrix},\qquad
 v=\begin{pmatrix}w\\\alpha^2\end{pmatrix},\quad w=(1,\alpha)\trans.
\]
The leading $2\times2$ matrix $D$ is positive definite. The kernel equation
forces $b=-Dw/\alpha^2$ and $c=w\trans Dw/\alpha^4$. Thus
\begin{equation}\label{eq:localPSD}
 H_i=
 \begin{pmatrix}I_2\\-w\trans/\alpha^2\end{pmatrix}
 D
 \begin{pmatrix}I_2&-w/\alpha^2\end{pmatrix}.
\end{equation}
This is positive semidefinite of rank two, with exactly the stated kernel.
The exact certificates verify $H_iv=0$ separately as well.
\end{proof}

\begin{table}[ht]
\centering
\caption{Certified strict lower bounds; every displayed decimal is an exact
rational number, not a floating-point estimate. Indices inside $H_i$ are
one-based.}\label{tab:local}
\begin{tabular}{ccc}
\toprule
$i$ & Lower bound for $(H_i)_{11}$ & Lower bound for
$\det(H_i)_{\{1,2\},\{1,2\}}$\\
\midrule
1&1.016353&0.444291\\
2&1.150668&0.687760\\
3&0.380276&0.104352\\
4&0.242452&0.030000\\
5&3.486402&4.967538\\
6&0.424643&0.091575\\
7&1.071944&0.579814\\
\bottomrule
\end{tabular}
\end{table}

\section{A rational cone with seven irrational zero rays}\label{sec:cone}
Define the exact rationals
\begin{equation}\label{eq:triangleconstants}
 a_0=\frac{125992104989487316477}{10^{20}},\qquad
 b_0=\frac{158740105196819947475}{10^{20}},\qquad
 \delta=\frac1{10000},
\end{equation}
and the rational matrix
\begin{equation}\label{eq:W}
 W=\begin{pmatrix}
 1&1&1\\
 a_0+\delta&a_0&a_0-\delta\\
 b_0&b_0+\delta&b_0-\delta
 \end{pmatrix}.
\end{equation}
Its determinant is $3\delta^2>0$. Let $w_1,w_2,w_3$ be its columns.
For each $i$, put
\begin{equation}\label{eq:GiK}
 g_{ik}=C_iw_k\in\Q^7,\qquad
 G_i=\cone\{g_{i1},g_{i2},g_{i3}\},\qquad
 K=G_1+\cdots+G_7.
\end{equation}
Equivalently, $K$ is the cone generated by the $21$ columns of
\begin{equation}\label{eq:V}
 V=[C_1W\mid C_2W\mid\cdots\mid C_7W]\in\Q^{7\times21}.
\end{equation}
These generators need not be nonnegative in their seven-dimensional
coordinates. Their nonnegative embedding is obtained from the facet matrix
$R$ in the next section.

\subsection{Exact inclusion and cross-pairing certificates}
Put
\[
 d_x=\frac{\alpha-a_0}{\delta},\qquad
 d_y=\frac{\alpha^2-b_0}{\delta},
\]
and define
\begin{equation}\label{eq:barycentric}
 \lambda_1=\frac{1+2d_x-d_y}{3},\qquad
 \lambda_2=\frac{1-d_x+2d_y}{3},\qquad
 \lambda_3=\frac{1-d_x-d_y}{3}.
\end{equation}
Direct calculation gives
\begin{equation}\label{eq:inclusion}
 \lambda_1+\lambda_2+\lambda_3=1,\qquad
 W\lambda=v,\qquad u_i=\sum_{k=1}^3\lambda_kg_{ik}.
\end{equation}
The rational interval certificate gives $\lambda_k>33/100$ for all three
indices. Hence $u_i\in\ri G_i$.

All cross-block generator pairings satisfy the following stronger-than-needed
strict inequality:
\begin{equation}\label{eq:cross}
 g_{ik}\trans Qg_{j\ell}>\frac{17}{1000}
 \quad(1\leq i<j\leq7,\ 1\leq k,\ell\leq3).
\end{equation}
There are exactly $\binom72\,3^2=189$ such checks. They are evaluated from
$V\trans QV$ in the same exact field arithmetic, using rational lower
endpoints; the smallest certified lower endpoint exceeds $0.017344$.

Finally, the rational vector
\begin{equation}\label{eq:h}
 h=\frac1{10^8}
 \begin{pmatrix}
 -71246841\\-101280383\\-150296097\\-110993458\\
 -39802169\\-119030527\\-63286873
 \end{pmatrix}
\end{equation}
satisfies
\begin{equation}\label{eq:pointedcertificate}
 h\trans g_{ik}>\frac{999}{1000}>0
 \qquad(1\leq i\leq7,\ 1\leq k\leq3).
\end{equation}
This last check is purely rational. These exact values, the barycentric
coefficients, and the generator matrix are supplied in
\texttt{data/rational\_cone\_certificate.json}.

\begin{proposition}[The zero set on the rational cone]\label{prop:zeros}
The cone $K$ is pointed, rational polyhedral, and full-dimensional in $\R^7$.
The quadratic form $q(x)=x\trans Qx$ is nonnegative on $K$, and
\begin{equation}\label{eq:zeros}
 \{x\in K:q(x)=0\}
 =\{0\}\ \cup\ \bigcup_{i=1}^7\{t u_i:t>0\}.
\end{equation}
In particular, its only rational zero in $K$ is the origin.
\end{proposition}

\begin{proof}
Finite rational generation makes $K$ rational polyhedral. The common
strictly positive functional \eqref{eq:pointedcertificate} makes it pointed:
any nonzero conic combination of generators has positive $h$-value, so
$K$ cannot contain both a nonzero vector and its negative. By
\eqref{eq:inclusion}, $K$ contains all columns of the orthogonal matrix $O$.
They span $\R^7$, proving full dimension.

For $x_i\in G_i\subset U_i$, write $x_i=C_iz_i$. Lemma~\ref{lem:seed} gives
\[
 q(x_i)=z_i\trans H_iz_i\geq0,
\]
with equality precisely when $x_i\in\R u_i$. Because $C_i$ is injective
and every generator $w_k$ has first coordinate one, the first coordinate
of any nonzero vector in $\cone\{w_1,w_2,w_3\}$ is strictly positive.
Thus the only nonzero part of $\R u_i$ inside $G_i$ is its positive ray.

For nonzero $x_i\in G_i$ and $x_j\in G_j$ with $i\ne j$, write each as a
nonnegative combination of its three generators. Each combination has at
least one positive coefficient. Expanding the bilinear pairing and using
\eqref{eq:cross} gives
\begin{equation}\label{eq:crossGi}
 x_i\trans Qx_j>0.
\end{equation}
Now any $x\in K$ has some representation $x=x_1+\cdots+x_7$ with
$x_i\in G_i$, and
\begin{equation}\label{eq:qsum}
 q(x)=\sum_i q(x_i)+2\sum_{i<j}x_i\trans Qx_j\geq0.
\end{equation}
If at least two $x_i$ are nonzero, the right-hand side is strictly positive.
If exactly one is nonzero, it vanishes precisely on that group's positive
$u_i$-ray. The zero vector causes no exception because the common positive
functional excludes a nontrivial conic sum equal to zero. This proves
\eqref{eq:zeros}.

Each $C_i$ has rational column rank three. Lemma~\ref{lem:irrational}
shows that none of these seven lines contains a nonzero rational vector.
Consequently $K$ has no nonzero rational zero of $q$.
\end{proof}

\section{The exact facet matrix}\label{sec:facets}
The passage from a rational cone in $\R^7$ to an ordinary nonnegative Gram
problem must use a complete inequality description of the cone. Merely
finding some supporting inequalities would not suffice. The following
construction and exhaustive check supply that completeness.

Multiply each of the $21$ columns of $V$ by a positive rational number to
obtain a primitive integer vector. Denote the resulting vectors, in the
same block order, by $p_1,\ldots,p_{21}\in\Z^7$. No column changes its ray.
For every six-element subset $J\subset\{1,\ldots,21\}$, form the
$6\times7$ matrix with rows $p_j\trans$, $j\in J$.

If the row rank is six, compute a primitive integer generator $r_J$ of its
one-dimensional nullspace. Retain it precisely when one orientation has
\begin{equation}\label{eq:facetcriterion}
 r_J\trans p_k\geq0\qquad(1\leq k\leq21).
\end{equation}
Use that orientation and discard duplicate normals. Sort the retained
normals first by their increasing zero-index lists
$\{k:r_J\trans p_k=0\}$ and then lexicographically by the normal entries.
The transposes of these normals are the rows of $R$.

\begin{lemma}[Complete facet enumeration]\label{lem:facets}
This exact enumeration produces $444$ distinct primitive integer normals.
Every corresponding zero-index list has size six, and
\begin{equation}\label{eq:Hrepresentation}
 K=\{x\in\R^7:Rx\geq0\},\qquad
 R\in\Z^{444\times7},\quad \rank R=7.
\end{equation}
\end{lemma}

\begin{proof}
The finite enumeration gives the counts in Table~\ref{tab:facets}.
Nullspaces, ranks, orientations, and all supporting signs are calculated
using exact integers. In particular, the $444$ retained normals vanish on
six independent generators, so each defines a six-dimensional face of
$K$: a facet.

Conversely, every facet of a full-dimensional finitely generated cone in
$\R^7$ is generated by the input generators lying on that facet, and has
linear span of dimension six. It therefore contains six independent input
generators. Their index set occurs in the exhaustive enumeration, and
their oriented primitive null normal is the facet normal. No facet can be
missing.

A full-dimensional pointed polyhedral cone is the intersection of its
facet halfspaces. One can also see this here by slicing with $h\trans x=1$:
\eqref{eq:pointedcertificate} gives a compact six-dimensional polytope whose
vertices are among $p_k/(h\trans p_k)$. Its facet inequalities completely
determine the slice, and homogenizing gives the cone description. Thus
\eqref{eq:Hrepresentation} follows. The exact rational row-rank check gives
$\rank R=7$; it is also forced by pointedness of the intersection.
\end{proof}

\begin{table}[ht]
\centering
\caption{Exhaustive exact facet certificate.}\label{tab:facets}
\begin{tabular}{lr}
\toprule
Quantity & Verified value\\
\midrule
Six-generator subsets examined, $\binom{21}{6}$ & 54,264\\
Rank-deficient six-generator subsets & 0\\
Non-supporting six-generator subsets & 53,820\\
Distinct supporting primitive facet normals & 444\\
Generators on each facet & 6\\
Rank of the complete facet matrix & 7\\
\bottomrule
\end{tabular}
\end{table}

The file \texttt{data/facets.json} records both the oriented primitive
normals and their exact zero-index lists. The verifier does not trust a
stored ``complete'' flag: it independently recomputes all $54,264$
nullspaces and compares the resulting complete set with the stored one.
The construction of $R$ involves no numerical convex-hull oracle.

\section{Proof of the counterexample and further consequences}
\begin{proof}[Proof of Theorem~\ref{thm:main}]
Apply Lemma~\ref{lem:obstruction} with $r=7$, $N=444$, the cone from
\eqref{eq:GiK}, the orthogonal matrix \eqref{eq:cayley}, and the form
\eqref{eq:Q}. Lemma~\ref{lem:seed} gives $\tr Q=0$,
Proposition~\ref{prop:zeros} gives the required absence of nonzero rational
zeros, and Lemma~\ref{lem:facets} gives the exact rational cone embedding.
Therefore $A=RR\trans$ is completely positive, has no rational nonnegative
factor of any finite width, and lies on the boundary because $444>7$.

It remains to verify the claimed strict positivity and cp-rank. Let
$B=RO$. Every row of $R$ is nonnegative on every generator $g_{ik}$.
Since all three coefficients in \eqref{eq:inclusion} are positive,
\begin{equation}\label{eq:Bzeros}
 (Ru_i)_a=0
 \quad\Longleftrightarrow\quad
 \text{row $a$ of $R$ vanishes on all three }g_{i1},g_{i2},g_{i3}.
\end{equation}
A facet contains exactly six of the $21$ generator rays. Hence a row of
$B$ can vanish in at most two of its seven coordinates. Every row of $B$
therefore has at least five strictly positive coordinates. Any two rows
have a common strictly positive coordinate, since two subsets of size at
least five in a seven-element set intersect. Their inner product is
strictly positive. Thus every entry of $A=BB\trans=RR\trans$ is positive;
it is an integer because $R$ is an integer matrix.

Finally, $\rank A=\rank R=7$. The displayed nonnegative factor $B=RO$ has
seven columns, so $\cpr(A)\leq7$. Every Gram factor has at least
$\rank A$ columns, giving equality.
\end{proof}

\subsection{Why adding columns cannot repair rationality}
The obstruction actually determines all possible factor directions, not
just a selected factorization. If $A=CC\trans$ is any real nonnegative
factorization, the same reduction $X=LC$ gives $XX\trans=I_7$ and
$x_j\in K$. The trace calculation forces every nonzero $x_j$ onto one of
the seven rays in \eqref{eq:zeros}. Thus every nonzero column of $C$ lies
on one of
\begin{equation}\label{eq:ambientrays}
 \R_{>0}Ru_1,\ldots,\R_{>0}Ru_7.
\end{equation}
Each ambient ray also has no nonzero rational vector: applying the
rational left inverse $L$ to such a vector would give a nonzero rational
vector on $\R u_i$. Splitting a column into more columns along its ray
cannot avoid this obstruction.

In reduced coordinates, orthogonality makes the remaining freedom explicit.
Writing each nonzero $x_j=t_j u_{i(j)}$ with $t_j>0$, the identity
$XX\trans=I_7$ is equivalent to
\[
 \sum_{j:i(j)=i}t_j^2=1\qquad(1\leq i\leq7).
\]
There is freedom to split each unit weight, but none of the permitted
nonzero directions is rational.

\begin{corollary}[Larger orders]\label{cor:padding}
For every $n\geq444$, there is an entrywise strictly positive integer
matrix of order $n$ and rank seven that is completely positive but has no
rational nonnegative factor of any finite width.
\end{corollary}
\begin{proof}
Append copies of one row of $R$ to obtain $R_n\in\Z^{n\times7}$ and set
$A_n=R_nR_n\trans$. The factor $R_nO$ is nonnegative. All its row inner
products are positive by the preceding support argument. A rational
nonnegative factor of $A_n$ would restrict on its first $444$ rows to one
of $A$, which is impossible.
\end{proof}

\begin{corollary}[Nonclosedness of rational factorability]
For order $444$, the set of matrices with a finite nonnegative rational
Gram factor is not closed even relative to
$\CP_{444}\cap\Q^{444\times444}$ with its Euclidean subspace topology.
\end{corollary}
\begin{proof}
Approximate every entry of the finite nonnegative real factor $RO$ by
nonnegative rational numbers. The Gram matrices of the resulting rational
factors converge to $A$, but Theorem~\ref{thm:main} excludes $A$ itself.
\end{proof}

\section{Exact arithmetic, finite verification, and reproducibility}
\label{sec:verification}
The role of computation is restricted to explicit finite equalities,
ranks, signs, and an exhaustive facet list. The quantified conclusion
about all rational factors is the argument in
Lemma~\ref{lem:obstruction}; no enumeration of candidate factor widths
or rational matrices is involved.

\subsection{Cubic field arithmetic and sign certificates}
Represent an element of $\F$ by a triple of rational numbers
$(a,b,c)$ meaning $a+b\alpha+c\alpha^2$. Addition is coordinatewise,
and multiplication is the exact rule
\begin{align}\label{eq:fieldmultiply}
 &(a+b\alpha+c\alpha^2)(d+e\alpha+f\alpha^2)\nonumber\\
 &\quad=(ad+2bf+2ce)+(ae+bd+2cf)\alpha+(af+be+cd)\alpha^2.
\end{align}
Every rational is stored as an arbitrary-precision integer numerator and
positive integer denominator. Gaussian elimination uses exact inverses
in this cubic field; no algebraic-number approximations enter an equality
or rank test.

For inequalities, the verifier uses
\begin{equation}\label{eq:isolation}
 \begin{aligned}
 \ell&=\frac{12599210498948731647672106072782283505702}{10^{40}},\\[0.4em]
 u&=\frac{12599210498948731647672106072782283505703}{10^{40}}.
 \end{aligned}
\end{equation}
It checks $\ell^3<2<u^3$ by integer arithmetic. The real cube function is
strictly increasing, so $\ell<\alpha<u$. Each quadratic expression
$a+b\alpha+c\alpha^2$ is evaluated by interval Horner arithmetic on
$[\ell,u]$. Products use the minimum and maximum of the four endpoint
products, and all endpoints remain rational. A positivity assertion is
accepted only if its exact lower endpoint is positive.

The same algorithm certifies the principal minors of $H_i$, all $21$
off-diagonal entries of $Q'$, all $189$ cross-generator pairings, and the
barycentric coefficients. Dependency overestimation in interval
arithmetic can only make the bounds wider; it cannot invalidate a
positive lower endpoint.

\subsection{Facet arithmetic}
For each $6\times7$ integer matrix, the nullspace routine uses
fraction-free Bareiss elimination with row and column pivoting. Every
integer division is checked for zero remainder. Integer back substitution
produces a null vector, which is checked against all six input rows and
normalized by the positive gcd of its entries. Orienting and testing it
against all $21$ generators uses integer dot products only.

Both the original exact symbolic construction and the independent
standard-library reconstruction produce identical algebraic certificate
entries. The final distributed verifier requires no third-party packages:
it rebuilds the seed from \eqref{eq:S}, \eqref{eq:T}, and
\eqref{eq:freevalues}, rechecks the cone, and exhaustively reconstructs the
facet list. The mathematical proof does not depend on the numerical search
that first suggested the seed.

\subsection{Complete verification command}
From the extracted package directory, run
\begin{Verbatim}[fontsize=\small]
python code/verify_all.py
python code/test_arithmetic.py
\end{Verbatim}
The first command verifies the seed, all cone properties, all $54,264$
facet candidates, and every entry of the expanded matrix. The second runs
$18$ arithmetic and corruption-detection unit tests. The recorded runs
passed. Python's optimization switches \texttt{-O} and \texttt{-OO} are
explicitly rejected so that assertion checks cannot be disabled silently.

The verified matrix has $98,790$ stored lower-triangular entries, each a
strictly positive integer. The largest entry has $270$ decimal digits.
The factor $RO$ has $2,966$ strictly positive entries and $142$ exact zeros.
The Gram identity is certified by $OO\trans=I_7$ and the exact definition
$A=RR\trans$, and the expanded matrix file is checked entry by entry
against that definition.

The file \texttt{logs/verification\_summary.json} records exact rational
bounds and the verified counts. Unit-test results are in
\texttt{logs/unit\_tests.txt}. These logs are not accepted as inputs to
establish the theorem: running the verifier recomputes the checks.

\subsection{What is and is not being claimed}
The construction refutes the universal assertion in
\eqref{eq:question}. It does not assert that order $444$ or rank seven is
minimal. It does not classify every matrix of order five or the intermediate
ranks. Those additional classification questions are not necessary to
answer the stated universal existence question negatively.

The counterexample is singular but entrywise strictly positive. Singularity
is allowed by the stated boundary problem, and the unrestricted-width
condition has been addressed directly. The factor field
$\Q(\sqrt[3]{2})$ is not totally real: its two other embeddings are
nonreal. A descent argument that requires total reality therefore cannot
be applied to this factor. No field-descent theorem is used here.

The finite calculations are exact and reproducible. The complete
mathematical argument is presented above, but neither independent human
review nor formal proof-assistant verification has occurred. No priority
claim relative to all existing literature is made, and no repository
status or external file has been changed by this package.

\appendix
\section{A compact verification of the geometric argument}
The following implication chain is useful for auditing the interface
between the exact data and the unrestricted-width conclusion.

First, the arithmetic seed supplies seven rational three-dimensional
subspaces $U_i$ and an orthonormal basis $u_i\in U_i$. The form $Q$ is
positive semidefinite on each $U_i$, with kernel precisely $\R u_i$.
The column rank-three certificates imply that every such line is
irrational in the sense of Lemma~\ref{lem:irrational}.

Second, the triangle $W$ supplies rational cones $G_i\subset U_i$ containing
$u_i$ with strictly positive coefficients. The $189$ strict inequalities
force positive bilinear pairings between nonzero vectors in different
cones. Hence a zero of $Q$ in their sum cannot use more than one group.
This is the only step needed to pass from local semidefiniteness to the
complete zero-set description.

Third, the common rational functional proves pointedness and the
orthonormal basis proves full dimension. Exhaustive facet enumeration
then identifies the sum cone \emph{exactly} with $\{x:Rx\geq0\}$.
The inclusion of the basis gives $RO\geq0$, while the equality of the two
cone descriptions constrains \emph{every} possible reduced factor column,
not just the displayed seven-column factor.

Finally, rationality is preserved by the left inverse of the rational
full-column-rank matrix $R$. The trace-zero identity turns the nonnegative
quadratic values of every reduced factor column into individual zeros.
The irrational-ray test then rules out every nonzero rational column,
independently of how many columns there are.

None of these implications uses density of rational points to obtain an
exact boundary factor, a conjectural rational exposing normal, or an
assumption that rational cp-rank equals real cp-rank.

\begin{thebibliography}{9}
\bibitem{repo}
OpenProblemsInNLA,
\emph{PF-03---Rational factors for rational completely positive boundary
matrices}, repository statement, status check dated 10 September 2026.
\url{https://raw.githubusercontent.com/ajt60gaibb/OpenProblemsInNLA/main/nonnegative-and-positive-factorizations/PF-03/README.md}.

\bibitem{bsm}
A.~Berman and N.~Shaked-Monderer,
\emph{Completely Positive Matrices over Sets},
Communications in Optimization Theory (2025), article 23, Section 2.
\url{https://cot.mathres.org/issues/COT202523.pdf}.

\bibitem{dsv}
M.~Dutour Sikiri\'c, A.~Sch\"urmann, and F.~Vallentin,
\emph{Rational factorizations of completely positive matrices},
Linear Algebra and its Applications \textbf{523} (2017), 46--51.
\url{https://doi.org/10.1016/j.laa.2017.02.017}.
\end{thebibliography}
\end{document}
